\paragraph{Smoke-test trajectory.} We validate the entire pipeline on a
synthetic 4-turn trajectory in which an LLM cognitive map drifts on
turn~3 (chair coordinates corrupted) and approximately recovers on
turn~4. This is sufficient to exercise all metric paths and to verify
both directions of the \tb update machine
(Algorithm~\ref{alg:tb})---reinforce, buffer, overwrite. The
pipeline produces \verb|comparison_metrics.json| and
\verb|comparison_chart.png| (Figure~\ref{fig:metric}).

\begin{figure}[t]
\centering
\includegraphics[width=0.96\textwidth]{fig01_metric_comparison.pdf}
\caption{Smoke-test comparison on the synthetic trajectory. Five
\tos probing metrics show \tb-favouring deltas; one
(\textsc{Perception (Pos)}) goes the other way, reflecting the
expected sub-cell-precision cost of ego$\to$global re-projection.
\textbf{Real-data version follows on Day~2.}}\label{fig:metric}
\end{figure}

\begin{table}[t]
\centering\small
\caption{Smoke-test aggregate (single trajectory, $n{=}1$); positive
$\Delta$ means \tb is better.}\label{tab:smoke}
\begin{tabular}{lrrr}
\toprule
Metric & LLM & \tb & $\Delta$ \\\midrule
Correctness (Pos) & 0.26 & 0.44 & \textbf{+0.19} \\
Correctness (Ori) & 0.67 & 1.00 & \textbf{+0.33} \\
Stability (Pos)   & 0.67 & 1.00 & \textbf{+0.33} \\
Stability (Ori)   & 0.71 & 1.00 & \textbf{+0.29} \\
Perception (Pos)  & 0.84 & 0.59 & $-$0.25 \\
Perception (Ori)  & 1.00 & 1.00 & 0.00 \\
Self-track (Pos)  & 1.00 & 1.00 & 0.00 \\
Self-track (Ori)  & 1.00 & 1.00 & 0.00 \\
\bottomrule
\end{tabular}
\end{table}

The single negative $\Delta$ is the expected ego$\to$global
discretisation cost discussed in Section~\ref{sec:method} (sub-cell
precision is sacrificed in exchange for region-level stability and
orientation correctness). The four positive deltas on Correctness and
Stability are consistent with Proposition~\ref{prop:stab} for
$b{=}2$ at the configured $\epsilon$.

\begin{figure}[t]
\centering
\begin{subfigure}{0.46\textwidth}\centering
\includegraphics[width=\linewidth]{fig02_update_rule.pdf}
\caption{The four-state object-update machine of \tb.}
\end{subfigure}\hfill
\begin{subfigure}{0.46\textwidth}\centering
\includegraphics[width=\linewidth]{fig03_drift_curve_cartoon.pdf}
\caption{Hypothesised drift profile (cartoon, replaced by real
AutoDL data on Day~2).}
\end{subfigure}
\caption{Update rule (left) and predicted Stability shape (right).}
\end{figure}

\paragraph{What the smoke test does \emph{not} prove.}
The synthetic trajectory uses a "perfect-perception LLM" that copies
ground-truth coordinates verbatim, which is why
\textsc{Perception (Pos)} hits 0.84 for the LLM in
Table~\ref{tab:smoke}. On real AutoDL-rendered data, the LLM will
itself be subject to perception error, and the gap on
\textsc{Perception (Pos)} is expected to shrink or invert. The
preliminary numbers therefore validate \emph{the pipeline and the
direction of the comparisons}, not the magnitudes.
